% !TEX root = ../main.tex
%% Chapitre 15 — Intégration
%% Programme officiel MPSI/MP2I (BO spécial n°1 du 11/02/2021), deuxième semestre,
%% section « Intégration » : a) continuité uniforme, théorème de Heine ;
%% b) fonctions en escalier, continues par morceaux ; c) intégrale sur un segment
%% (valeurs dans K), propriétés ; d) sommes de Riemann (démonstration exigible
%% pour f lipschitzienne) ; e) lien intégrale-primitive ; f) formules de Taylor
%% globales (reste intégral, inégalité de Taylor-Lagrange).
%% Résultats vérifiés par verifications/chapitre_15.py.

\begin{ouverture}[integrale]{15}{Intégration}{Partie 2 · Deuxième semestre}
  {Continuité uniforme\enspace·\enspace Intégrale sur un segment\enspace·\enspace
   Sommes de Riemann\enspace·\enspace Primitives\enspace·\enspace Formules de Taylor}
  \objectif{Comprendre la continuité uniforme et le théorème de Heine.}
  \objectif{Construire l'intégrale des fonctions continues par morceaux sur un segment.}
  \objectif{Manier linéarité, positivité, croissance, Chasles et l'inégalité triangulaire.}
  \objectif{Reconnaître et calculer des limites de sommes de Riemann.}
  \objectif{Relier intégrale et primitives : le théorème fondamental de l'analyse.}
  \objectif{Appliquer Taylor reste intégral et l'inégalité de Taylor-Lagrange.}
  \prerequis{Techniques de calcul intégral : primitives, IPP, changement de variable (chapitre 4).}
  \prerequis{Suites et séries d'encadrements ; Bolzano-Weierstrass, Cesàro (chapitre 5).}
  \prerequis{Analyse asymptotique : équivalents, développements (chapitre 11).}
\end{ouverture}

%% =====================================================================
%%  L'ESSENTIEL
%% =====================================================================
\essentiel

\essSection{1}{Continuité uniforme}

\begin{definition}[Continuité uniforme]
$f$ est \textbf{uniformément continue} sur $I$ si
\[
\forall\varepsilon>0,\ \exists\delta>0,\ \forall x,y\in I,\qquad
|x-y|\leqslant\delta\implies|f(x)-f(y)|\leqslant\varepsilon :
\]
le même $\delta$ sert \emph{partout} — alors que la continuité autorise un
$\delta$ différent en chaque point. Toute fonction lipschitzienne est
uniformément continue ($\delta=\varepsilon/K$), toute fonction uniformément
continue est continue.
\end{definition}

\begin{theoreme}[Heine]
Toute fonction continue sur un \emph{segment} $\intFF{a}{b}$ y est
uniformément continue (démonstration non exigible — l'exercice 21 la
reconstruit avec Bolzano-Weierstrass). C'est la clé de la construction de
l'intégrale : sur un segment, une fonction continue se laisse approcher
\emph{uniformément} par des fonctions en escalier.
\end{theoreme}

\begin{exemple}
$\sqrt{\phantom{x}}$ est uniformément continue sur $\R_{+}$
($|\sqrt x-\sqrt y|\leqslant\sqrt{|x-y|}$) sans être lipschitzienne près de
$0$ ; $x\mapsto x^{2}$ est continue mais pas uniformément continue sur
$\R$ : les points $n$ et $n+\frac1n$ se rapprochent, leurs images s'écartent
(exercice 9).
\end{exemple}

\essSection{2}{Fonctions en escalier, continues par morceaux}

\begin{definition}[Subdivisions, escalier, morceaux]
Une \textbf{subdivision} de $\intFF{a}{b}$ est une suite finie
$a=x_0<x_1<\dots<x_n=b$ ; son \emph{pas} est le plus grand des
$x_{k+1}-x_k$. Une fonction est \textbf{en escalier} si elle est constante
sur chaque $\intOO{x_k}{x_{k+1}}$ d'une subdivision adaptée ;
\textbf{continue par morceaux} si sur chaque $\intOO{x_k}{x_{k+1}}$ elle
est continue et prolongeable par continuité aux bornes. Les fonctions
considérées sont définies sur un segment, à valeurs dans $\K=\R$ ou $\Cx$ ;
les fonctions continues par morceaux forment un sous-espace vectoriel et un
sous-anneau de $\K^{\intFF{a}{b}}$.
\end{definition}

\essSection{3}{L'intégrale sur un segment et ses propriétés}

\noindent
\begin{minipage}[c]{0.55\linewidth}
\begin{definition}[Intégrale d'une fonction continue par morceaux]
L'intégrale d'une fonction en escalier est la somme des aires algébriques
de ses rectangles ; celle d'une fonction continue par morceaux $f$ est la
valeur commune vers laquelle se resserrent les intégrales des fonctions en
escalier qui encadrent $f$ (le programme n'impose pas de construction
particulière). Notations :
$\displaystyle\int_{\intFF{a}{b}}f$, $\displaystyle\int_{a}^{b}f$,
$\displaystyle\int_{a}^{b}f(t)\,\mathrm{d}t$ — et, si $b\leqslant a$,
$\displaystyle\int_{a}^{b}f=-\int_{b}^{a}f$.
\end{definition}
\end{minipage}\hfill
\begin{minipage}[c]{0.41\linewidth}
\centering
\begin{tikzpicture}[scale=0.62,line cap=round]
  \fill[otGreen,opacity=0.3]
    plot[domain=0.5:4.4,samples=60] ({\x},{1.3+0.7*sin(deg(1.5*\x-1))+0.12*\x})
    -- (4.4,0) -- (0.5,0) -- cycle;
  \draw[otMuted,line width=0.5pt,-{Stealth[length=1.6mm]}] (-0.3,0) -- (5.2,0);
  \draw[otMuted,line width=0.5pt,-{Stealth[length=1.6mm]}] (0,-0.3) -- (0,3.1);
  \draw[otNavy,line width=1.2pt]
    plot[domain=0.15:4.9,samples=80] ({\x},{1.3+0.7*sin(deg(1.5*\x-1))+0.12*\x});
  \draw[otGreenDark,line width=1pt] (0.5,-0.12) -- (0.5,0.12);
  \draw[otGreenDark,line width=1pt] (4.4,-0.12) -- (4.4,0.12);
  \node[below,font=\scriptsize,text=otGreenDark] at (0.5,-0.12) {$a$};
  \node[below,font=\scriptsize,text=otGreenDark] at (4.4,-0.12) {$b$};
  \node[font=\scriptsize,text=otGreenDark!80!black] at (2.45,0.85) {$\displaystyle\int_a^b f$};
\end{tikzpicture}\\[0.5mm]
\legende{L'intégrale est l'aire \emph{algébrique} entre la courbe et
l'axe : comptée négativement sous l'axe.}
\end{minipage}

\begin{theoreme}[Les propriétés fondamentales]
Pour $f,g$ continues par morceaux sur $\intFF{a}{b}$ ($a\leqslant b$) :
\begin{itemize}[label={\color{otNavy}\textbullet},leftmargin=1.4em,itemsep=0.6mm,topsep=0.8mm]
  \item \emph{linéarité} :
    $\int(\lambda f+\mu g)=\lambda\int f+\mu\int g$ ;
  \item \emph{positivité et croissance} : si $f\geqslant0$ alors
    $\int f\geqslant0$ ; si $f\leqslant g$ alors
    $\int f\leqslant\int g$ ;
  \item \emph{inégalité triangulaire} :
    $\Bigl|\displaystyle\int_{a}^{b}f\Bigr|\leqslant\displaystyle\int_{a}^{b}|f|$ ;
  \item \emph{relation de Chasles} :
    $\displaystyle\int_{a}^{c}f=\int_{a}^{b}f+\int_{b}^{c}f$ (valable pour
    toutes positions de $a,b,c$ avec la convention d'orientation) ;
  \item si $f$ est \emph{continue}, positive et d'intégrale nulle, alors
    $f=0$ (faux pour une fonction continue par morceaux : un point isolé
    ne pèse rien).
\end{itemize}
La \textbf{valeur moyenne} de $f$ est
$\dfrac{1}{b-a}\displaystyle\int_{a}^{b}f$. Enfin l'intégrale d'une
fonction impaire sur un segment centré en $0$ est nulle, celle d'une paire
vaut deux fois l'intégrale sur la moitié, et l'intégrale d'une fonction
périodique sur un intervalle de période ne dépend pas de l'intervalle.
\end{theoreme}

\essSection{4}{Sommes de Riemann}

\noindent
\begin{minipage}[c]{0.55\linewidth}
\begin{theoreme}[Sommes de Riemann]
Si $f$ est continue par morceaux sur $\intFF{a}{b}$ :
\[
\frac{b-a}{n}\sum_{k=0}^{n-1}f\Bigl(a+k\,\frac{b-a}{n}\Bigr)
\xrightarrow[n\to+\infty]{}\int_{a}^{b}f(t)\,\mathrm{d}t .
\]
(On peut aussi évaluer aux points droits ou en tout point de chaque
maille.) La démonstration est \emph{exigible} pour $f$ lipschitzienne —
la voici.
\end{theoreme}
\end{minipage}\hfill
\begin{minipage}[c]{0.41\linewidth}
\centering
\begin{tikzpicture}[scale=0.62,line cap=round]
  \foreach \k in {0,...,7}{%
    \pgfmathsetmacro{\xa}{0.5+0.4875*\k}%
    \pgfmathsetmacro{\hh}{1.3+0.7*sin(deg(1.5*\xa-1))+0.12*\xa}%
    \fill[otGreen,opacity=0.28] ({\xa},0) rectangle ({\xa+0.4875},{\hh});
    \draw[otGreenDark,line width=0.5pt] ({\xa},0) rectangle ({\xa+0.4875},{\hh});}
  \draw[otMuted,line width=0.5pt,-{Stealth[length=1.6mm]}] (-0.3,0) -- (5.2,0);
  \draw[otMuted,line width=0.5pt,-{Stealth[length=1.6mm]}] (0,-0.3) -- (0,3.1);
  \draw[otNavy,line width=1.2pt]
    plot[domain=0.15:4.9,samples=80] ({\x},{1.3+0.7*sin(deg(1.5*\x-1))+0.12*\x});
  \node[below,font=\scriptsize,text=otGreenDark] at (0.5,-0.12) {$a$};
  \node[below,font=\scriptsize,text=otGreenDark] at (4.4,-0.12) {$b$};
\end{tikzpicture}\\[0.5mm]
\legende{La méthode des rectangles : quand le pas tend vers $0$, l'aire
des rectangles tend vers l'intégrale.}
\end{minipage}

\begin{demonstration}[Cas lipschitzien — démonstration exigible]
Soit $K$ une constante de Lipschitz de $f$, $h=\frac{b-a}{n}$ et
$x_k=a+kh$. Sur chaque maille,
\[
\Bigl|\int_{x_k}^{x_{k+1}}f(t)\,\mathrm{d}t-h\,f(x_k)\Bigr|
=\Bigl|\int_{x_k}^{x_{k+1}}\bigl(f(t)-f(x_k)\bigr)\mathrm{d}t\Bigr|
\leqslant\int_{x_k}^{x_{k+1}}K\,(t-x_k)\,\mathrm{d}t=\frac{K h^{2}}{2}.
\]
En sommant les $n$ mailles (Chasles) :
$\Bigl|\int_{a}^{b}f-S_n\Bigr|\leqslant n\cdot\frac{Kh^{2}}{2}
=\frac{K(b-a)^{2}}{2n}\xrightarrow[n\to+\infty]{}0$.
\end{demonstration}

\begin{savoirfaire}[Reconnaître une somme de Riemann]
Mettre $\dfrac1n$ en facteur et faire apparaître $f\bigl(\tfrac kn\bigr)$ :
\[
\sum_{k=1}^{n}\frac{1}{n+k}
=\frac1n\sum_{k=1}^{n}\frac{1}{1+\frac kn}
\xrightarrow[n\to+\infty]{}\int_{0}^{1}\frac{\mathrm{d}t}{1+t}=\ln2 .
\]
Les produits se traitent en passant au logarithme (exercice 11).
\end{savoirfaire}

\essSection{5}{Intégrale et primitives}

\begin{theoreme}[Théorème fondamental de l'analyse]
Soit $f$ \emph{continue} sur un intervalle $I$ et $a\in I$. La fonction
$F:x\mapsto\displaystyle\int_{a}^{x}f(t)\,\mathrm{d}t$ est de classe
$\mathcal{C}^{1}$ sur $I$, avec $F'=f$ : \emph{toute fonction continue sur
un intervalle possède des primitives}. En conséquence, pour toute primitive
$G$ de $f$ :
\[
\int_{a}^{b}f(t)\,\mathrm{d}t=G(b)-G(a) ,
\]
et les théorèmes d'intégration par parties et de changement de variable du
chapitre 4 reposent sur cette égalité. Si la borne est composée,
$x\mapsto\int_{a}^{u(x)}f$ se dérive en $u'(x)\,f\bigl(u(x)\bigr)$.
\end{theoreme}

\begin{attention}
L'hypothèse de continuité est essentielle pour $F'=f$ : pour $f$ seulement
continue par morceaux, $F$ est continue et dérivable sauf aux points de
discontinuité. Et une intégrale à valeurs complexes s'intègre par parties
réelle et imaginaire : $\int\mathrm{e}^{\mathrm{i}t}\mathrm{d}t$ se calcule
comme en réel, ce qui unifie bien des calculs trigonométriques
(exercice 18).
\end{attention}

\essSection{6}{Formules de Taylor globales}

\begin{theoreme}[Taylor avec reste intégral ; inégalité de Taylor-Lagrange]
Soit $f$ de classe $\mathcal{C}^{n+1}$ sur $I$, et $a,b\in I$ :
\[
f(b)=\sum_{k=0}^{n}\frac{f^{(k)}(a)}{k!}(b-a)^{k}
+\int_{a}^{b}\frac{(b-t)^{n}}{n!}\,f^{(n+1)}(t)\,\mathrm{d}t .
\]
En majorant le reste : si $\bigl|f^{(n+1)}\bigr|\leqslant M$ entre $a$ et
$b$,
\[
\Bigl|f(b)-\sum_{k=0}^{n}\frac{f^{(k)}(a)}{k!}(b-a)^{k}\Bigr|
\leqslant M\,\frac{|b-a|^{n+1}}{(n+1)!}
\qquad\text{(inégalité de Taylor-Lagrange).}
\]
\end{theoreme}

\begin{attention}
Taylor-Young (chapitre 11) est un énoncé \emph{local} : il décrit $f$ près
de $a$, sans contrôle de l'erreur. Les formules globales ci-dessus donnent
une \emph{majoration effective} valable sur tout l'intervalle : ce sont
elles qui prouvent $\mathrm{e}=\sum1/k!$, encadrent les restes et
justifient les calculs approchés (exercices 7, 8, 15). L'égalité de
Taylor-Lagrange est hors programme.
\end{attention}

\begin{savoirfaire}[Comparer les méthodes d'approximation]
Pour approcher $\int_a^b f$ avec $n$ mailles : les \emph{rectangles} font
une erreur en $O\bigl(\frac1n\bigr)$ (la borne lipschitzienne ci-dessus
est optimale en général), les \emph{trapèzes} — moyenne des deux rectangles
— une erreur en $O\bigl(\frac{1}{n^{2}}\bigr)$ pour $f$ assez régulière :
à coût égal, les trapèzes gagnent largement (exercice 16).
\end{savoirfaire}

%% =====================================================================
%%  EXERCICES
%% =====================================================================
\exercices
\rubrique{Applications}

\begin{exercice}[1]{Fonctions en escalier}
\begin{questions}
  \item Calculer $\displaystyle\int_{0}^{3}\lfloor t\rfloor\,\mathrm{d}t$
    et $\displaystyle\int_{0}^{3}\bigl(t-\lfloor t\rfloor\bigr)\mathrm{d}t$.
  \item Les fonctions suivantes sont-elles continues par morceaux sur
    $\intFF{0}{1}$ : $x\mapsto|2x-1|$ ; $x\mapsto\lfloor3x\rfloor$ ;
    $x\mapsto\dfrac1x$ prolongée par $0$ en $0$ ?
\end{questions}
\end{exercice}

\begin{exercice}[1]{Les propriétés de base}
\begin{questions}
  \item Sans calculer l'intégrale, encadrer
    $\displaystyle\int_{0}^{1}\frac{\mathrm{d}x}{1+x}$ entre $\dfrac12$ et
    $1$ ; comparer à sa valeur exacte.
  \item Justifier
    $\Bigl|\displaystyle\int_{0}^{1}\sin\bigl(x^{2}\bigr)\mathrm{d}x\Bigr|\leqslant1$.
  \item Calculer $\displaystyle\int_{-1}^{2}|x|\,\mathrm{d}x$ avec la
    relation de Chasles.
\end{questions}
\end{exercice}

\begin{exercice}[1]{Parité et périodicité}
\begin{questions}
  \item Que vaut $\displaystyle\int_{-1}^{1}x^{3}\cos x\,\mathrm{d}x$ ?
  \item Calculer $\displaystyle\int_{0}^{2\pi}\cos^{2}t\,\mathrm{d}t$.
  \item Montrer que $\displaystyle\int_{a}^{a+2\pi}\sin t\,\mathrm{d}t$ ne
    dépend pas de $a$ — et donner sa valeur.
\end{questions}
\end{exercice}

\begin{exercice}[1]{Valeur moyenne}
\begin{questions}
  \item Calculer la valeur moyenne de $\sin$ sur $\intFF{0}{\pi}$, puis
    celle de $x\mapsto x^{2}$ sur $\intFF{0}{1}$.
  \item Interpréter géométriquement : quel rectangle a la même aire que la
    région sous la courbe ?
\end{questions}
\end{exercice}

\begin{exercice}[1]{Sommes de Riemann : premiers calculs}
Déterminer les limites des suites :
\begin{tasks}(3)
  \task $\displaystyle\sum_{k=1}^{n}\frac{1}{n+k}$
  \task $\displaystyle\sum_{k=1}^{n}\frac{k}{n^{2}}$
  \task $\displaystyle\sum_{k=1}^{n}\frac{n}{n^{2}+k^{2}}$
\end{tasks}
\end{exercice}

\begin{exercice}[1]{Dériver une intégrale}
\begin{questions}
  \item Soit $F(x)=\displaystyle\int_{0}^{x}\mathrm{e}^{-t^{2}}\mathrm{d}t$.
    Justifier que $F$ est de classe $\mathcal{C}^{1}$, donner $F'$ et les
    variations de $F$.
  \item Soit $G(x)=\displaystyle\int_{0}^{x^{2}}\sin t\,\mathrm{d}t$.
    Calculer $G'(x)$.
\end{questions}
\end{exercice}

\begin{exercice}[1]{Taylor avec reste intégral}
\begin{questions}
  \item Écrire la formule de Taylor avec reste intégral pour $\exp$ entre
    $0$ et $1$ à l'ordre $n$.
  \item En déduire
    $0<\mathrm{e}-\displaystyle\sum_{k=0}^{n}\frac{1}{k!}
    \leqslant\frac{\mathrm{e}}{(n+1)!}$, et le nombre de termes assurant
    une précision de $10^{-3}$.
\end{questions}
\end{exercice}

\begin{exercice}[1]{L'inégalité de Taylor-Lagrange}
\begin{questions}
  \item Montrer que $|\sin x-x|\leqslant\dfrac{|x|^{3}}{6}$ pour tout
    $x\in\R$, et que
    $\Bigl|\cos x-1+\dfrac{x^{2}}{2}\Bigr|\leqslant\dfrac{x^{4}}{24}$.
  \item En déduire une valeur approchée de $\sin(0{,}1)$ et la précision
    garantie.
\end{questions}
\end{exercice}

\rubrique{Entraînement}

\begin{exercice}[2]{Continuité uniforme : exemples et contre-exemples}
\begin{questions}
  \item Montrer que $|\sqrt x-\sqrt y|\leqslant\sqrt{|x-y|}$ pour
    $x,y\geqslant0$, et en déduire que la racine carrée est uniformément
    continue sur $\R_{+}$. Est-elle lipschitzienne sur $\intFF{0}{1}$ ?
  \item Montrer que $x\mapsto x^{2}$ n'est pas uniformément continue sur
    $\R$. \emph{Considérer $x_n=n$ et $y_n=n+\frac1n$.}
  \item Montrer que $x\mapsto\sin\dfrac1x$ n'est pas uniformément continue
    sur $\intOF{0}{1}$.
\end{questions}
\end{exercice}

\begin{exercice}[2]{Intégrale nulle}
\begin{questions}
  \item Soit $f$ continue et positive sur $\intFF{a}{b}$ avec
    $\displaystyle\int_{a}^{b}f=0$. Montrer que $f=0$.
    \emph{Si $f(c)>0$, minorer $f$ sur un petit segment autour de $c$.}
  \item En déduire : si $f$ est continue et
    $\displaystyle\int_{a}^{b}f^{2}=0$, alors $f=0$.
  \item Donner une fonction continue par morceaux, positive, non nulle,
    d'intégrale nulle.
\end{questions}
\end{exercice}

\begin{exercice}[2]{Sommes de Riemann corsées}
Déterminer les limites :
\begin{questions}
  \item $\displaystyle\frac{1}{n\sqrt n}\sum_{k=1}^{n}\sqrt{k}$ ;
  \item $\displaystyle\Bigl(\prod_{k=1}^{n}\Bigl(1+\frac kn\Bigr)\Bigr)^{1/n}$
    \emph{(passer au logarithme)} ;
  \item $\displaystyle\frac1n\sum_{k=1}^{n}\sin\frac{k\pi}{n}$.
\end{questions}
\end{exercice}

\begin{exercice}[2]{La démonstration lipschitzienne, en pratique}
Soit $f$ lipschitzienne de constante $K$ sur $\intFF{a}{b}$ et $S_n$ la
somme de Riemann aux points gauches.
\begin{questions}
  \item Redémontrer l'estimation
    $\Bigl|\displaystyle\int_{a}^{b}f-S_n\Bigr|
    \leqslant\dfrac{K(b-a)^{2}}{2n}$.
  \item L'appliquer à $\displaystyle\int_{0}^{1}\mathrm{e}^{x}\mathrm{d}x$ :
    quelle valeur de $n$ garantit une erreur inférieure à $10^{-2}$ ?
    Comparer à l'erreur réellement observée.
\end{questions}
\end{exercice}

\begin{exercice}[2]{La suite $I_n=\int_0^1 x^{n}\mathrm{e}^{x}\,\mathrm{d}x$}
\begin{questions}
  \item Établir la relation $I_n=\mathrm{e}-n\,I_{n-1}$ pour
    $n\geqslant1$.
  \item Encadrer $I_n$ entre $\dfrac{1}{n+1}$ et $\dfrac{\mathrm{e}}{n+1}$,
    et en déduire sa limite.
  \item Montrer que $n\,I_{n-1}\to\mathrm{e}$, puis que
    $I_n\sim\dfrac{\mathrm{e}}{n}$ (chapitre 11).
\end{questions}
\end{exercice}

\begin{exercice}[2]{Autour de $\int_0^1\frac{\mathrm{d}x}{1+x^{n}}$}
On pose $K_n=\displaystyle\int_{0}^{1}\frac{\mathrm{d}x}{1+x^{n}}$ et
$J_n=1-K_n=\displaystyle\int_{0}^{1}\frac{x^{n}}{1+x^{n}}\,\mathrm{d}x$.
\begin{questions}
  \item Montrer que $K_n\to1$. \emph{Encadrer $J_n$.}
  \item Par le changement de variable $u=x^{n}$, montrer que
    $0\leqslant\ln2-nJ_n\leqslant\dfrac{1}{n+1}$.
  \item En déduire le développement
    $K_n=1-\dfrac{\ln2}{n}+o\Bigl(\dfrac1n\Bigr)$.
\end{questions}
\end{exercice}

\begin{exercice}[2]{$\mathrm{e}$ et $\ln2$ : restes et irrationalité}
\begin{questions}
  \item Avec l'exercice 7, montrer que
    $A_n=n!\,\displaystyle\sum_{k=0}^{n}\frac{1}{k!}$ est un entier
    vérifiant $0<n!\,\mathrm{e}-A_n<1$ pour $n\geqslant2$, et en déduire
    que $\mathrm{e}$ est irrationnel (chapitre 5, revisité).
  \item Montrer que le reste de la série alternée de $\ln2$ est exactement
    l'intégrale $\displaystyle\int_{0}^{1}\frac{x^{n}}{1+x}\,\mathrm{d}x$
    du chapitre 11 — Taylor reste intégral appliqué à $\ln(1+x)$.
\end{questions}
\end{exercice}

\begin{exercice}[2]{Rectangles contre trapèzes}
Pour $\displaystyle\int_{0}^{1}\mathrm{e}^{x}\mathrm{d}x=\mathrm{e}-1$, on
note $R_n$ la méthode des rectangles (points gauches) et
$T_n$ celle des trapèzes.
\begin{questions}
  \item Exprimer $T_n$ à partir des sommes aux points gauches et droits.
  \item Calculer numériquement les erreurs pour $n=10$ puis $n=100$, et
    vérifier qu'elles se divisent par environ $10$ pour $R_n$, par environ
    $100$ pour $T_n$ : ordres $1$ et $2$.
\end{questions}
\end{exercice}

\begin{exercice}[2]{L'inégalité de Cauchy-Schwarz intégrale}
Pour $f,g$ continues sur $\intFF{a}{b}$ :
\begin{questions}
  \item En étudiant le trinôme
    $t\mapsto\displaystyle\int_{a}^{b}(f+tg)^{2}$, montrer que
    $\Bigl(\displaystyle\int_{a}^{b}fg\Bigr)^{2}
    \leqslant\displaystyle\int_{a}^{b}f^{2}\int_{a}^{b}g^{2}$.
  \item En déduire
    $\Bigl(\displaystyle\int_{0}^{1}f\Bigr)^{2}
    \leqslant\displaystyle\int_{0}^{1}f^{2}$ et interpréter en termes de
    moyennes.
  \item Quand y a-t-il égalité ? (On y reviendra au chapitre 18.)
\end{questions}
\end{exercice}

\begin{exercice}[2]{Fonctions à valeurs complexes}
\begin{questions}
  \item Calculer $\displaystyle\int_{0}^{\pi}\mathrm{e}^{\mathrm{i}t}\,\mathrm{d}t$
    et vérifier l'inégalité
    $\bigl|\int\bigr|\leqslant\int|\cdot|$ sur cet exemple.
  \item Retrouver une primitive de $x\mapsto\mathrm{e}^{x}\cos x$ en
    intégrant $x\mapsto\mathrm{e}^{(1+\mathrm{i})x}$ (chapitres 3 et 4).
\end{questions}
\end{exercice}

\rubrique{Approfondissement}

\begin{exercice}[3]{La constante d'Euler, enfin}
On pose $u_n=H_n-\ln n$, où $H_n=1+\frac12+\dots+\frac1n$.
\begin{questions}
  \item Par comparaison de $\dfrac1k$ aux intégrales
    $\displaystyle\int_{k}^{k+1}\frac{\mathrm{d}t}{t}$, montrer que
    $(u_n)$ est décroissante et minorée par $0$.
  \item En déduire la convergence de $(u_n)$ : sa limite est la constante
    d'Euler $\gamma$, admise au chapitre 5 — désormais construite.
  \item Montrer l'encadrement $\gamma<u_n<\gamma+\dfrac1n$ et commenter la
    vitesse de convergence (chapitre 11).
\end{questions}
\end{exercice}

\begin{exercice}[3]{Le lemme de Riemann-Lebesgue, version $\mathcal{C}^{1}$}
Soit $f$ de classe $\mathcal{C}^{1}$ sur $\intFF{a}{b}$.
\begin{questions}
  \item Par intégration par parties, majorer
    $\Bigl|\displaystyle\int_{a}^{b}f(t)\sin(nt)\,\mathrm{d}t\Bigr|$ par
    une constante divisée par $n$, et conclure que l'intégrale tend vers
    $0$.
  \item Calculer exactement
    $\displaystyle\int_{0}^{2\pi}t^{2}\sin(nt)\,\mathrm{d}t$ et vérifier.
  \item Pourquoi le résultat vaut-il aussi pour les fonctions en escalier,
    puis continues par morceaux ?
\end{questions}
\end{exercice}

\begin{exercice}[3]{Le théorème de Heine, démonstration guidée}
Soit $f$ continue sur $\intFF{a}{b}$, supposée \emph{non} uniformément
continue.
\begin{questions}
  \item Montrer qu'il existe $\varepsilon>0$ et deux suites $(x_n)$,
    $(y_n)$ de $\intFF{a}{b}$ telles que $|x_n-y_n|\leqslant\frac1n$ et
    $|f(x_n)-f(y_n)|>\varepsilon$.
  \item Avec Bolzano-Weierstrass (chapitre 5), extraire une sous-suite
    convergente et aboutir à une contradiction.
\end{questions}
\end{exercice}

\begin{exercice}[3]{Orthogonale aux polynômes}
Soit $f$ continue sur $\intFF{0}{1}$ telle que
$\displaystyle\int_{0}^{1}f(t)\,t^{k}\,\mathrm{d}t=0$ pour
$k=0,1,\dots,n$.
\begin{questions}
  \item Montrer que $\displaystyle\int_{0}^{1}fP=0$ pour tout
    $P\in\R_n[X]$.
  \item Montrer que $f$ s'annule en au moins $n+1$ points distincts de
    $\intOO{0}{1}$. \emph{Sinon, construire un polynôme de degré
    $\leqslant n$ ayant partout le signe de $f$, et contredire
    l'exercice 10.}
  \item Vérifier sur $f(x)=6x^{2}-6x+1$ (orthogonale à $1$ et $x$) :
    combien de racines dans $\intOO{0}{1}$ ?
\end{questions}
\end{exercice}

\begin{exercice}[3]{L'aire de l'ellipse}
\begin{questions}
  \item Paramétrer l'ellipse $\dfrac{x^{2}}{a^{2}}+\dfrac{y^{2}}{b^{2}}=1$
    par $x=a\cos t$, $y=b\sin t$, et exprimer son aire comme l'intégrale
    $\displaystyle\int_{0}^{2\pi}ab\sin^{2}t\,\mathrm{d}t$.
  \item Conclure : aire $=\pi ab$, et retrouver le point de vue
    « déterminant » du chapitre 14.
\end{questions}
\end{exercice}

\begin{exercice}[3]{La formule de Stirling, enfin démontrée}
On pose $u_n=\dfrac{n!}{\sqrt n\,(n/\mathrm{e})^{n}}$ et
$d_n=\ln u_n-\ln u_{n+1}$.
\begin{questions}
  \item Montrer que $d_n=\Bigl(n+\dfrac12\Bigr)\ln\Bigl(1+\dfrac1n\Bigr)-1$.
  \item Montrer que $\ln(1+t)\geqslant\dfrac{2t}{2+t}$ pour $t\geqslant0$,
    et en déduire $d_n\geqslant0$ ; montrer par ailleurs
    $d_n\leqslant\dfrac{1}{6n^{2}}$ pour $n \geqslant 2$
    \emph{(série alternée de $\ln(1+x)$)}.
  \item En déduire que $(\ln u_n)$ converge
    \emph{(comparer $\sum d_n$ à $\sum\frac{1}{n^{2}}$, chapitre 5)}, puis
    que $u_n\to\ell>0$ : il existe une constante $\ell$ telle que
    $n!\sim\ell\,\sqrt n\,\bigl(n/\mathrm{e}\bigr)^{n}$.
  \item À l'aide de l'équivalent de Wallis
    $\binom{2n}{n}\sim\dfrac{4^{n}}{\sqrt{\pi n}}$ (chapitre 11, obtenu
    sans Stirling), montrer que $\ell=\sqrt{2\pi}$ : la formule de
    Stirling, admise au chapitre 11, est démontrée.
\end{questions}
\end{exercice}

% !TEX root = ../../main.tex
%% Défi du chapitre 15 — L'intégrale de Gauss
%% Résultats vérifiés par verifications/defis_11_15.py.

\rubrique{Défi}

\begin{defi}{L'intégrale de Gauss}{2 h}
La fonction $t\mapsto\mathrm{e}^{-t^{2}}$ n'a pas de primitive exprimable à
l'aide des fonctions usuelles (théorème de Liouville, admis) ; on va
pourtant calculer exactement l'aire totale sous sa courbe, la plus célèbre
des probabilités. On pose
\[
F(A)=\int_0^{A}\mathrm{e}^{-t^{2}}\,\mathrm{d}t\quad(A\geqslant0),
\qquad
W_k=\int_0^{\pi/2}\cos^{k}\theta\,\mathrm{d}\theta\quad(k\in\N),
\]
ces dernières étant les intégrales de Wallis, dont on rappelle
l'équivalent $W_k\sim\sqrt{\pi/(2k)}$ (chapitre 11, exercice 23).

\partieExo{A}{Existence et encadrement}
\begin{questions}
  \item Montrer que $F$ est croissante et que $F(A)\leqslant1+\frac1{\mathrm{e}}$
    pour tout $A\geqslant0$. \emph{Comparer $\mathrm{e}^{-t^{2}}$ et
    $\mathrm{e}^{-t}$ sur $\intFO{1}{+\infty}$.} En déduire que $F(A)$ admet
    une limite finie $G$ quand $A\to+\infty$.
  \item Montrer que $1+u\leqslant\mathrm{e}^{u}$ pour tout réel $u$. En
    déduire, pour $n\geqslant1$ et $t\in\intFF{0}{\sqrt n}$ :
    \[
    \Bigl(1-\frac{t^{2}}{n}\Bigr)^{n}\leqslant\mathrm{e}^{-t^{2}}
    \leqslant\Bigl(1+\frac{t^{2}}{n}\Bigr)^{-n} .
    \]
\end{questions}

\medskip
\begin{center}
\begin{tikzpicture}[x=2.3cm,y=2.6cm,line cap=round]
  \fill[otGreen,opacity=0.18] (0,0) -- plot[domain=0:2.6,samples=90] ({\x},{exp(-\x*\x)}) -- (2.6,0) -- cycle;
  \draw[otMuted,line width=0.5pt,-{Stealth[length=1.6mm]}] (0,0) -- (2.8,0);
  \draw[otMuted,line width=0.5pt,-{Stealth[length=1.6mm]}] (0,0) -- (0,1.2);
  \draw[otNavy!60,line width=0.9pt,dash pattern=on 3pt off 1.5pt] plot[domain=0:2.6,samples=90]
    ({\x},{(1+\x*\x/3)^(-3)});
  \draw[otNavy,line width=0.9pt,dash pattern=on 1pt off 1.2pt] plot[domain=0:1.7320,samples=90]
    ({\x},{(1-\x*\x/3)^3});
  \draw[otGreenDark,line width=1.3pt] plot[domain=0:2.6,samples=90] ({\x},{exp(-\x*\x)});
  \draw[otMuted,line width=0.4pt] (1.7320,0.03) -- (1.7320,-0.03);
  \node[font=\tiny,text=otMuted,below] at (1.7320,-0.02) {$\sqrt3$};
  \node[font=\tiny,text=otMuted,anchor=east] at (0,1) {$1$};
  \node[font=\scriptsize,text=otGreenDark,anchor=south west] at (0.62,0.72) {$\mathrm{e}^{-t^{2}}$};
  \node[font=\scriptsize,text=otNavy!70,anchor=west] at (1.25,0.33) {$\bigl(1+\frac{t^{2}}{3}\bigr)^{-3}$};
  \node[font=\scriptsize,text=otNavy,anchor=east] at (0.5,0.28) {$\bigl(1-\frac{t^{2}}{3}\bigr)^{3}$};
  \draw[otNavy,line width=0.4pt] (0.5,0.28) -- (0.86,0.31);
\end{tikzpicture}\\[0.5mm]
\legende{La courbe de Gauss (en vert, aire $G$ ombrée) encadrée par les
deux fonctions de la question 2 pour $n=3$.}
\end{center}

\partieExo{B}{Le calcul par les intégrales de Wallis}
\begin{questions}[resume]
  \item Par le changement de variable $t=\sqrt n\sin\theta$, montrer que
    \[
    \int_0^{\sqrt n}\Bigl(1-\frac{t^{2}}{n}\Bigr)^{n}\mathrm{d}t
    =\sqrt n\,W_{2n+1} .
    \]
  \item Par le changement de variable $t=\sqrt n\tan\theta$, montrer que,
    pour $n\geqslant2$,
    \[
    \int_0^{\sqrt n}\Bigl(1+\frac{t^{2}}{n}\Bigr)^{-n}\mathrm{d}t
    \leqslant\sqrt n\,W_{2n-2} .
    \]
  \item En déduire $\sqrt n\,W_{2n+1}\leqslant F\bigl(\sqrt n\bigr)\leqslant\sqrt n\,W_{2n-2}$,
    puis
    \[
    G=\lim_{A\to+\infty}\int_0^{A}\mathrm{e}^{-t^{2}}\,\mathrm{d}t=\frac{\sqrt\pi}{2} .
    \]
\end{questions}

\partieExo{C}{Applications}
\begin{questions}[resume]
  \item Pour $a>0$, déterminer la limite de
    $\displaystyle\int_0^{A}\mathrm{e}^{-at^{2}}\,\mathrm{d}t$ quand
    $A\to+\infty$.
  \item Montrer que
    $\dfrac{1}{\sqrt{2\pi}}\displaystyle\int_{-A}^{A}\mathrm{e}^{-t^{2}/2}\,\mathrm{d}t
    \xrightarrow[A\to+\infty]{}1$ : c'est la normalisation de la \emph{loi
    normale}, limite universelle des sommes de variables aléatoires
    indépendantes (programme de deuxième année).
  \item D'où vient le $\sqrt\pi$ ? Comparer avec les apparitions de $\pi$
    dans la formule de Stirling (chapitre 15) et dans la marche aléatoire
    (chapitre 17).
\end{questions}
\end{defi}

